% ============================================================
%  ch7_VI  末章：用初等变换构造满足截断条件的迭代 $\varphi$-格
%  原书 pp. 46--56，即 PDF 第 47--57 页；转录稿 partF.md
%  说明：本章自第 46 页上某个证明的末尾接续，转录稿开头没有章标题，
%        故本文件不设 \section，直接从正文写起。
% ============================================================

\begin{proof}[证明]
\origpage{46}与退化子\footnote{译注：本文件自第 46 页开头接续，上一页的末句在此页续完（原文跨页，句子未完），故开头不是一个完整的句子；该句所属证明的陈述与前半也不在本段转录稿内。原段末尾的 $\square$ 即该证明的结束符，此处由 proof 环境自动生成。}；并且注意到，只要我们为族中每个类型 $\underline{r}$ 的某个迭代 $\varphi$-格 $M$ 都装上了水平结构 $N$，则所有迭代 $\varphi$-格就都装上了，因为总可以通过沿各 $\widehat{W}_w$ 的初等变换从一个过渡到另一个。

命题的结论于是来自下述事实：不存在无穷长的迭代 $\varphi$-格链
\[
M = M_0 \subsetneq M_1 \subsetneq \dots \subsetneq M_k \subsetneq \dots
\]
其中对每个指标 $k$，$M_{k+1}$ 都是 $M_k$ 经 $\widehat{W}_{w_0}$ 的反向变换，并且嵌入
\[
\mathcal{E}^{M_k}_{w_0} \hookrightarrow \mathcal{E}^{M_{k+1}}_{w_0}, \qquad
\mathcal{E}^{M_{k+1}}/\mathcal{E}^{M_{k+1}}_{w_0} \hookrightarrow \mathcal{E}^{M_k}/\mathcal{E}^{M_k}_{w_0}
\]
都是同构。
\end{proof}

有了前三个命题，就可以把迭代 $\varphi$-格沿正向或反向的若干串初等变换组合起来，得到：

\begin{proposition*}
设 $M$ 是 $V$ 中的一个迭代 $\varphi$-格，$\{\widehat{W}_w\}$ 是 $V$ 的一个由子 $\varphi$-空间给出的滤过，它是 $M$ 的典范滤过与一个元素 $\widehat{W}_{w_0}$ 的并，且该元素满足
\[
\deg \mathcal{E}^M_{w_0} > d(w_0).
\]
设对 $\{\widehat{W}_w\}$ 的秩列中每个 $w > w_0$（相应地，每个 $w \neq w_0$）都有
\[
\overline{\mathcal{E}}^M_w + {}^\tau\mathcal{E}^M_w \subsetneq {}^\tau\mathcal{E}^M.
\]
于是，在把 $A$ 换成一个只依赖于滤过 $\{\widehat{W}_w\}$ 以及各分次对象 $\mathcal{E}^M_{w^+}/\mathcal{E}^M_{w^-}$ 的最小斜率 $\mu^-$ 与最大斜率 $\mu^+$ 的有限扩张之后，可以从 $M$ 出发经一串初等变换构造出迭代 $\varphi$-格 $M'$，使得：

\begin{itemize}
\item $M'$ 所对应的 $\widehat{V}$ 的滤过包含于 $\{\widehat{W}_w\}$，并且含有元素 $\widehat{W}_{w_0}$。

\origpage{47}\item 对每个使 $\widehat{W}_w$ 出现在 $M'$ 所对应的 $\widehat{V}$ 的滤过中的秩 $w > w_0$（相应地，$w \neq w_0$），都有等式 $\deg \mathcal{E}^{M'}_w = \deg \mathcal{E}^M_w$，以及同构
\[
\begin{cases}
\mathcal{E}^M/\mathcal{E}^M_w \xrightarrow{\ \sim\ } \mathcal{E}^{M'}/\mathcal{E}^{M'}_w, \quad
\overline{\mathcal{E}}^M_w \xrightarrow{\ \sim\ } \overline{\mathcal{E}}^{M'}_w, & \text{若 } w > w_0,\\[1mm]
\mathcal{E}^{M'}_w \xrightarrow{\ \sim\ } \mathcal{E}^M_w, \quad
{}^\tau\mathcal{E}^{M'}_w/\overline{\mathcal{E}}^{M'}_w \xrightarrow{\ \sim\ } {}^\tau\mathcal{E}^M_w/\overline{\mathcal{E}}^M_w, & \text{若 } w < w_0,
\end{cases}
\]
从而有
\[
\overline{\mathcal{E}}^{M'}_w + {}^\tau\mathcal{E}^{M'}_w \subsetneq {}^\tau\mathcal{E}^{M'}.
\]

\item 有 $\deg \mathcal{E}^{M'}_{w_0} = \deg \mathcal{E}^M_{w_0} - 1$，以及
\[
\overline{\mathcal{E}}^{M'}_{w_0} + {}^\tau\mathcal{E}^{M'}_{w_0} \subsetneq {}^\tau\mathcal{E}^{M'}.
\]

\item 有典范嵌入
\[
\mathcal{E}^M_{w^+}/\mathcal{E}^M_w \hookrightarrow \mathcal{E}^{M'}_{w^+}/\mathcal{E}^{M'}_w, \quad \forall w \ge w_0
\]
以及
\[
\mathcal{E}^{M'}_{w_0} \hookrightarrow \mathcal{E}^M_{w_0}
\]
（相应地，还有 $\mathcal{E}^{M'}_w/\mathcal{E}^{M'}_{w^-} \hookrightarrow \mathcal{E}^M_w/\mathcal{E}^M_{w^-}$，$\forall w \le w_0$）。
\end{itemize}
\end{proposition*}

\begin{proof}[证明]
先证第一个断言。

由于 $\deg \mathcal{E}^M_{w_0} > d(w_0)$，$M$ 经 $\widehat{W}_{w_0}$ 有严格变换 $M^1$。

若典范嵌入 $\mathcal{E}^M/\mathcal{E}^M_{w_0^+} \hookrightarrow \mathcal{E}^{M^1}/\mathcal{E}^{M^1}_{w_0^+}$ 的余核为 $0$，就令 $M^2 = M^1$。

反之，若它是 $\kappa_A$ 上的 $1$ 维空间，就记 $M^2$ 为 $M^1$ 经 $\widehat{W}_{w_0}$ 的反向变换（它存在——必要时像命题中那样把 $A$ 换成有限扩张——因为必有 $\overline{\mathcal{E}}^{M^1}_{w_0^+} + {}^\tau\mathcal{E}^{M^1}_{w_0^+} \subsetneq {}^\tau\mathcal{E}^{M^1}$）。

然后令 $w_3 = w_0^+$，$M^3$ 为 $M^2$ 在秩 $w_3$ 处的极大修正；再令 $w_4 = w_3^+$，$M^4$ 为 $M^3$ 在秩 $w_4$ 处的极大修正，如此继续，直到达到某个秩 $w_k$，使得 $w_k^+ = r$，或者 $\widehat{W}_{w_k}$ 已出现在 $M^k$ 所对应的 $\widehat{V}$ 的滤过中，此时停止。

注意到有典范同构
\[
\mathcal{E}^M/\mathcal{E}^M_{w_0^+} \xrightarrow{\ \sim\ } \mathcal{E}^{M^2}/\mathcal{E}^{M^2}_{w_0^+}.
\]

\origpage{48}若嵌入 $\mathcal{E}^M/\mathcal{E}^M_{w_0^+} \hookrightarrow \mathcal{E}^{M^1}/\mathcal{E}^{M^1}_{w_0^+}$ 是同构，则这是显然的，因为此时我们已令 $M^2 = M^1$。相反的情形，只需证明嵌入
\[
{}^\tau\mathcal{E}^{M^2}/{}^\tau\mathcal{E}^{M^2}_{w_0^+} \cong \overline{\mathcal{E}}^{M^1}_{w_0^+}
\hookrightarrow \overline{\mathcal{E}}^{M^1}_{w_0}/\overline{\mathcal{E}}^{M^1}_{w_0} \cap {}^\tau\mathcal{E}^{M^1}_{w_0^+}
\cong {}^\tau\mathcal{E}^M/{}^\tau\mathcal{E}^M_{w_0^+}
\]
是同构\footnote{译注：原文如此。同一式中 $\overline{\mathcal{E}}^{M^1}$ 的下标一处作 $w_0$、一处作 $w_0^+$，与并列出现的 ${}^\tau\mathcal{E}^{M^1}_{w_0^+}$ 也不一致；下文则统一用 $w_0^+$。此处照录原文。}。而我们有
\[
\overline{\mathcal{E}}^{M^1}_{w_0^+} \oplus {}^\tau\mathcal{E}^{M^1}_{w_0^+}
\subseteq \overline{\mathcal{E}}^{M^1}_{w_0} + {}^\tau\mathcal{E}^{M^1}_{w_0^+}
\subsetneq {}^\tau\mathcal{E}^{M^1}
\]
其中第一个包含关系其实是等号；由此得到 $\overline{\mathcal{E}}^{M^1}_{w_0^+} \cong \overline{\mathcal{E}}^{M^1}_{w_0}/\overline{\mathcal{E}}^{M^1}_{w_0} \cap {}^\tau\mathcal{E}^{M^1}_{w_0^+}$，因为商 ${}^\tau\mathcal{E}^{M^1}_{w_0^+}/(\overline{\mathcal{E}}^{M^1}_{w_0^+} \oplus {}^\tau\mathcal{E}^{M^1}_{w_0^+})$ 的维数至多为 $1$。

由此容易看出：

\begin{itemize}
\item 或者 $M' = M^k$ 满足命题中要求的全部条件；

\item 或者它除第三条外全部满足，此时有等式 $\deg \mathcal{E}^{M'}_{w_0} = \deg \mathcal{E}^M_{w_0}$ 以及典范同构
\[
\mathcal{E}^M_{w^+}/\mathcal{E}^M_w \xrightarrow{\ \sim\ } \mathcal{E}^{M'}_{w^+}/\mathcal{E}^{M'}_w, \quad \forall w \ge w_0,
\]
\[
\mathcal{E}^{M'}_{w_0} \xrightarrow{\ \sim\ } \mathcal{E}^M_{w_0}.
\]
\end{itemize}

第一种情形已经完成。第二种情形，从 $M'$ 出发重复我们从 $M$ 出发做过的事（无须再次扩大环 $A$），如此继续。由于 $A$ 不变，必然在有限步后终止。

现在来看第二个断言。

先对 $M$ 施行上述过程，得到一个满足第一个断言结论的迭代 $\varphi$-格 $M^1$。

若典范嵌入 $\mathcal{E}^{M^1}_{w_0^-} \longrightarrow \mathcal{E}^M_{w_0^-}$ 的余核为 $0$，就令 $M^2 = M^1$。

反之，若它是 $\kappa_A$ 上的 $1$ 维空间，就记 $M^2$ 为 $M^1$ 经 $\widehat{W}_{w_0^-}$ 的反向变换（它存在，并且满足 $\deg \mathcal{E}^{M^2}_{w_0^-} = \deg \mathcal{E}^{M^1}_{w_0^-} + 1 = \deg \mathcal{E}^M_{w_0^-}$，因为必有 $\overline{\mathcal{E}}^{M^1}_{w_0^-} + {}^\tau\mathcal{E}^{M^1}_{w_0^-} \subsetneq {}^\tau\mathcal{E}^{M^1}$）。

\origpage{49}若 $\deg \mathcal{E}^{M^2}_{w_0} = \deg \mathcal{E}^{M^1}_{w_0} = \deg \mathcal{E}^M_{w_0} - 1$，就令 $M^3$ 为 $M^2$ 在秩 $w_0$ 处的极大修正；反之，若 $\deg \mathcal{E}^{M^2}_{w_0} = \deg \mathcal{E}^{M^1}_{w_0} + 1 = \deg \mathcal{E}^M_{w_0}$，就令 $M^3$ 为 $M^2$ 经 $\widehat{W}_{w_0}$ 的严格变换。$M^3$ 所对应的 $\widehat{V}$ 的滤过含有元素 $\widehat{W}_{w_0}$（否则 $M^3$ 会由 $M$ 经一串反向初等变换得到，从而有 $\deg \mathcal{E}^{M^3}_{w_0} \geq \deg \mathcal{E}^M_{w_0}$）。

然后令 $w_4 = w_0^-$，$M^4$ 为 $M$ 在秩 $w_4$ 处的极大修正\footnote{译注：原文如此。按上下文，$M^4$ 应由 $M^3$ 而非 $M$ 得到，此处疑为 $M^3$。}；再令 $w_5 = w_4^-$，$M^5$ 为 $M^4$ 在秩 $w_5$ 处的极大修正，如此继续，直到达到某个秩 $w_k$，使得 $\widehat{W}_{w_k}$ 仍出现在 $M_k$ 所对应的 $\widehat{V}$ 的滤过中，或者 $w_k^- = 0$，此时停止。

迭代 $\varphi$-格 $M' = M^k$ 满足 $\deg \mathcal{E}^{M'}_{w_0} = \deg \mathcal{E}^{M^1}_{w_0} = \deg \mathcal{E}^M_{w_0} - 1$，并且处于下述两种情形之一：

\begin{itemize}
\item 或者，使 $\widehat{W}_w$ 出现在 $M'$ 所对应的 $\widehat{V}$ 的滤过中的最大秩 $w < w_0$ 满足 $\deg \mathcal{E}^{M'}_w = \deg \mathcal{E}^M_w$，此时 $M'$ 满足命题中要求的全部条件；

\item 或者，这个最大秩 $w < w_0$ 满足 $\deg \mathcal{E}^{M'}_w = \deg \mathcal{E}^M_w - 1$。
\end{itemize}

第一种情形已经完成。第二种情形，从 $M'$ 出发重复我们从 $M^1$ 出发做过的事，如此继续。由于环 $A$ 不变，必然在有限步后终止。
\end{proof}

应当注意，在这个命题的构造中，$\{\widehat{W}_w\}$ 中满足「$\widehat{W}_w$ 仍出现在 $M'$ 所对应的 $\widehat{V}$ 的滤过中」的那些秩 $w > w_0$（相应地，秩 $w < w_0$）构成的子集，恰由比某个秩更大（相应地，更小）的一切秩组成。

$\{\widehat{W}_w\}$ 中满足 $\deg \mathcal{E}^{M'}_w = \deg \mathcal{E}^M_w$ 的秩 $w > w_0$（相应地，秩 $w < w_0$）所构成的两个子集（它们包含上述子集）也一样。在所有其余的秩 $w$ 处（这些秩因而构成一个含 $w_0$ 的区间），有 $\deg \mathcal{E}^{M'}_w = \deg \mathcal{E}^M_w - 1$。

上述命题还有如下补充：

\begin{corollary*}
在命题的情形下，再设对 $\{\widehat{W}_w\}$ 的每个秩 $w > w_0$ 都有
\[
\overline{\mathcal{E}}^M_w + {}^\tau\mathcal{E}^M_{w^+} = {}^\tau\mathcal{E}^M.
\]
\origpage{50}于是，对每个使 $\widehat{W}_w$ 仍出现在 $M'$ 所对应的 $\widehat{V}$ 的滤过中的秩 $w > w_0$，始终有
\[
\overline{\mathcal{E}}^{M'}_w + {}^\tau\mathcal{E}^{M'}_{w^+} = {}^\tau\mathcal{E}^{M'}
\]
并且还获得了下述附加性质：
\[
\overline{\mathcal{E}}^{M'}_{w_0} + {}^\tau\mathcal{E}^{M'}_{w_0^+} = {}^\tau\mathcal{E}^{M'}.
\]
\end{corollary*}

\begin{proof}[证明]
对每个使 $\widehat{W}_w$ 出现在 $M'$ 所对应的滤过中的秩 $w > w_0$，都有典范同构
\[
\mathcal{E}^M/\mathcal{E}^M_w \xrightarrow{\ \sim\ } \mathcal{E}^{M'}/\mathcal{E}^{M'}_w, \qquad
\overline{\mathcal{E}}^M_w \xrightarrow{\ \sim\ } \overline{\mathcal{E}}^{M'}_w,
\]
因此只需证明
\[
\overline{\mathcal{E}}^{M'}_{w_0} + {}^\tau\mathcal{E}^{M'}_{w_0^+} = {}^\tau\mathcal{E}^{M'}.
\]

只需在命题第一个断言的框架内做这件事；为此参看 \cite{lafforgue98} 第 2c 段命题 12 证明的末尾。
\end{proof}

现在可以给出叠 $\mathrm{Cht}^{r,d,\vec{p}\le p}$ 的紧合性赋值判据之存在性部分的证明了。

回忆截断多边形
\[
p : [0, r] \to \mathbb{R}^+
\]
是凸的，而且按假设还满足
\[
[p(r') - p(r'-1)] - [p(r'+1) - p(r')] \geq 2, \qquad \forall r', \ 0 < r' < r.
\]

对 $\tilde{\mathcal{E}}$ 的一般纤维 $V$ 中每个迭代 $\varphi$-格 $M$，$\mathcal{E}(M)$ 模 $\pi_A$ 的归约 $\mathcal{E}^M$ 带有其典范滤过 $(\mathcal{E}^M_w)_{w \in \underline{r}^-_M \cup \underline{r}^+_M}$。对 $\mathcal{E}^M$ 的任一由“好”子对象构成、且加细前一滤过的滤过 $(\mathcal{E}^M_w)_{w \in \underline{w}^- \cup \underline{w}^+}$\footnote{译注：原文如此。前一处滤过的下标集为 $\underline{r}^-_M \cup \underline{r}^+_M$，此处却作 $\underline{w}^- \cup \underline{w}^+$；按上下文，二者似应统一。}，可以联系它的多边形 $[0, r] \to \mathbb{R}$，即连接下列各点的折线：
\[
w \mapsto \deg \mathcal{E}^M_w - \frac{w}{r}d, \qquad w \in \underline{w}^- \cup \underline{w}^+.
\]

\origpage{51}在所有这样形成的多边形中，有一个大于其余一切。记它为 $\overline{p}^M$，可以称之为 $M$ 的典范多边形。它不必是凸的，但在每个区间 $[r', r'^+]$（$r' \in \underline{r}^-_M$）上的限制是凸的，并且满足
\[
\overline{p}^M(w) = \deg \mathcal{E}^M_w - \frac{w}{r}d, \qquad \forall w \in \underline{r}^-_M \cup \underline{r}^+_M.
\]

可以用数
\[
\mu^M = \sum_{0 \le r' < r'' < r} \max\{0, [\overline{p}^M(r''+1) - \overline{p}^M(r'')] - [\overline{p}^M(r'+1) - \overline{p}^M(r')]\}
\]
来度量它偏离凸性的程度；这个数取值于 $\frac{1}{r}\mathbb{N}$，等于 $0$ 当且仅当多边形 $\overline{p}^M$ 是凸的。

本段开头列出了迭代 $\varphi$-格 $M$ 为定义 $\mathrm{Cht}^{r,d,\vec{p}\le p}(A)$ 中延拓 $\tilde{\mathcal{E}}$ 的点所必须满足的性质 (1)–(5)。

\begin{proof}[证明：这样的 $\varphi$-格的存在性]
先选取一个迭代 $\varphi$-格 $M^0$，其类型 $\underline{r}_{M^0}$ 极小，意思是：若 $M$ 是能与 $M^0$ 用一串正向或反向的初等变换连接起来的迭代 $\varphi$-格，则它的类型 $\underline{r}_M$ 不可能满足 $\underline{r}^+_M \subsetneq \underline{r}^+_{M^0}$。

必要时在所有秩 $w \in \underline{r}^-_{M^0} \cap \underline{r}^+_{M^0}$ 处把 $M^0$ 反向变换足够多次（由于 $\underline{r}_{M^0}$ 极小，这样的反向变换可以无限重复），便可设
\[
\deg \mathcal{E}^{M^0}_w \geq d(w) + 1, \qquad \forall w \in \underline{r}^-_{M^0} \cap \underline{r}^+_{M^0}.
\]

从 $M^0$ 出发，按下述方式递归构造一列迭代 $\varphi$-格 $M^0, M^1, \dots, M^k, \dots$：

设 $M^k$ 已构造好，考虑它的典范多边形 $\overline{p}^{M^k}$。

\begin{itemize}
\item 若 $\overline{p}^{M^k} \leq p$，就停止构造。

\item 否则，$\overline{p}^{M^k}(r') - p(r')$（$0 \leq r' \leq r$）的最大值是一个 $\ge 1$ 的整数。若以 $w_k$ 记取到该最大值的最大整数 $r'$，则在 $\mathcal{E}^{M^k}$ 中存在唯一的秩 $w_k$ 的“好”子对象 $\mathcal{E}^{M^k}_{w_k}$，使得
\[
\deg \mathcal{E}^{M^k}_{w_k} - \frac{w_k}{r}d = \overline{p}^{M^k}(w_k).
\]
\end{itemize}

\origpage{52}于是对 $M^k$ 及 $\mathcal{E}^{M^k}$ 的“好”子对象 $\mathcal{E}^{M^k}_{w_k}$ 施行前一命题中（第二个断言的）构造，定义迭代 $\varphi$-格 $M^{k+1}$，它特别满足

\begin{itemize}
\item 秩 $w_k$ 出现在 $M^{k+1}$ 的类型中，并且
\[
\deg \mathcal{E}^{M^{k+1}}_{w_k} = \deg \mathcal{E}^{M^k}_{w_k} - 1,
\]

\item 每个出现在 $M^{k+1}$ 的类型中的秩 $w \neq w_k$ 也出现在 $M^k$ 的类型中，并且
\[
\deg \mathcal{E}^{M^{k+1}}_w = \deg \mathcal{E}^{M^k}_w,
\]

\item 在每个出现在 $M^k$ 的类型中的秩 $w$ 处，有
\[
\deg \mathcal{E}^{M^{k+1}}_w =
\begin{cases}
\deg \mathcal{E}^{M^k}_w - 1, \\
\text{或} \quad \deg \mathcal{E}^{M^k}_w,
\end{cases}
\]

\item 在 $\mathcal{E}^{M^k}$ 与 $\mathcal{E}^{M^{k+1}}$ 的分次对象之间，沿一个方向或另一方向都有典范嵌入；除 $w_k$ 两侧各恰有一个（其余核为 $\kappa_A$ 上 $1$ 维的）之外，它们全是同构。
\end{itemize}

我们已经可以设 $M^0$ 满足性质 (1)；由前一命题，所有 $M^k$ 也都如此。它们还满足与 (1) 等价的性质 (2)。

由构造以及 Harder–Narasimhan 典范多边形的通常性质可知，所有 $M^k$ 都满足：

\begin{itemize}
\item 对每个属于 $M^k$ 的类型的秩 $w$，有
\[
\deg \mathcal{E}^{M^k}_w = \overline{p}^{M^k}(w) + \frac{w}{r}d \ge d(w).
\]

\item 总有
\[
\mu^{M^{k+1}} \le \mu^{M^k}
\]
并且若 $\mu^{M^{k+1}} = \mu^{M^k}$，则必有
\[
\overline{p}^{M^{k+1}} \le \overline{p}^{M^k} \quad \text{且} \quad \overline{p}^{M^{k+1}} \neq \overline{p}^{M^k}.
\]
\end{itemize}

\origpage{53}这就迫使 $M^k$ 这列在某个指标 $k_0$ 处停止。迭代 $\varphi$-格 $M^{k_0} = M'$ 于是回答了所提的问题：

已经知道它满足性质 (1) 与 (2)。

此外，它的典范多边形 $\overline{p}^{M'}$ 满足不等式
\[
\overline{p}^{M'} \le p
\]
因此 $M'$ 满足性质 (4)，$\overline{p}^{M'}$ 是凸多边形，并且对 $\mathcal{E}^{M'}$ 的每个“好”子对象 $\mathcal{E}^{M'}_w$ 都有
\[
\deg \mathcal{E}^{M'}_w \le d(w).
\]

剩下只需证明：若 $\mathcal{E}^{M'}_w$ 是类型 II 的“好”子对象，则甚至有
\[
\deg \mathcal{E}^{M'}_w \le d(w) - 1
\]
并且 $M'$ 满足性质 (3)。

考虑使
\[
\max_{0 \le r' \le r}\{\overline{p}^{M^k}(r') - p(r')\} = 1
\]
的指标 $k < k_0$ 的集合。

它们构成区间 $[k_1, k_0[$；并且所有多边形 $\overline{p}^{M^k}$（$k_1 \le k < k_0$）自动是凸的。

有
\[
\deg(\mathcal{E}^{M^{k_1}}_w) = d(w) + 1, \qquad \forall w \in \underline{r}^-_{M^{k_1}} \cap \underline{r}^+_{M^{k_1}},
\]
因此秩 $w_{k_1}$ 不小于 $M^{k_1}$ 的类型 $\underline{r}^-_{M^{k_1}} \cap \underline{r}^+_{M^{k_1}}$ 的所有秩；而秩列 $w_k$（$k_1 \le k < k_0$）严格递减。由前一推论，每个 $M^{k+1}$（$k_1 \le k < k_0$）在 $\underline{r}^-_{M^{k+1}} \cap \underline{r}^+_{M^{k+1}}$ 的所有秩 $w \ge w_k$ 处满足性质 (3)；最终 $M'$ 完整地满足性质 (3)。

最后，设 $\mathcal{E}^{M'}_w$ 是 $\mathcal{E}^{M'}$ 的“好”子对象，满足
\[
\deg \mathcal{E}^{M'}_w = d(w),
\]
并设 $\widehat{W}_w$ 是 $\widehat{V}$ 中与之相伴的子空间。由于
\[
\overline{p}^{M^k} \le p + 1, \qquad k_1 \le k \le k_0,
\]

\origpage{54}$\widehat{W}_w$ 与所有 $M^k$（$k_1 \le k \le k_0$）所对应的 $\widehat{V}$ 的滤过相容。

若 $w \le w_{k_0-1}$，“好”子对象 $\mathcal{E}^{M'}_w$ 必为类型 I。否则，存在整数 $k_2$（$k_1 \le k_2 < k_0$）使得
\[
w_{k_2} < w \le w_{k_2-1} \quad \text{若 } k_1 < k_2,
\]
或
\[
w_{k_1} < w \le r \quad \text{若 } k_2 = k_1,
\]
并且必有
\[
d(w) = \deg \mathcal{E}^{M'}_w \le \deg \mathcal{E}^{M^{k_0-1}}_w \le \dots \le \deg \mathcal{E}^{M^{k_2}}_w \le d(w),
\]
于是这些不等式全是等号。

这只有在“好”子对象 $\mathcal{E}^{M'}_w$ 为类型 I 时才可能。事实上，$M'$ 的类型 $\underline{r}^-_{M'}$ 中小于 $w$ 的最大秩，或者等于 $0$（此时 $\mathcal{E}^{M'}_w$ 为类型 I），或者形如 $w_{k_3}$，其中 $k_1 \le k_2 \le k_3 < k_0$。后一情形，${}^\tau\mathcal{E}^{M^{k_3}}/{}^\tau\mathcal{E}^{M^{k_3}}_w$ 与 $\overline{\mathcal{E}}^{M'}_{w_{k_3}}/\overline{\mathcal{E}}^{M'}_{w_{k_3}} \cap {}^\tau\mathcal{E}^{M'}_w$ 带有两个滤过，其分次对象同构；而度数等式
\[
\deg \mathcal{E}^{M^{k_3}}_w = \deg \mathcal{E}^{M'}_w
\]
迫使
\[
\overline{\mathcal{E}}^{M'}_{w_{k_3}} + {}^\tau\mathcal{E}^{M'}_w = {}^\tau\mathcal{E}^{M'},
\]
这意味着 $\mathcal{E}^{M'}_w$ 是类型 I。
\end{proof}

\begin{remark*}
\begin{itemize}
\item 可以看到，这里每个 $M^{k+1}$ 都从 $M^k$ 出发，选取使 $r' \mapsto \overline{p}^{M^k}(r') - p(r')$ 取到最大值的诸秩中最大者 $w_k$ 而构造。

这与 \cite{lafforgue98} 一文（第 1034 页延拓存在性证明的末尾）所说的程序不同：那里取的是 $\mathcal{E}^{M^k}$ 的一个满足 $\deg \mathcal{E}^{M^k}_w > d(w)$、且使该性质成立的秩 $w$ 最大的“好”子对象 $\mathcal{E}^{M^k}_w$。正如 Yakov Varshavsky 向作者指出的，这个程序一般行不通（从秩 $r = 3$ 起就不行）。必须在每一步都像我们刚才这样选取秩 $w_k$。

\origpage{55}\item \cite{lafforgue02} 一文中引理 V.20(iii) 的证明，其实质是下述\underline{具体}程序：从一个（非迭代的）$\varphi$-格 $M$ 出发，构造出定义 $\mathrm{Cht}^{r,d,\vec{p}\le p}(A)$ 中延拓的迭代 $\varphi$-格 $M'$；这里 $M$ 对应 $\mathrm{Cht}^{r,d}(A)$ 中的一个延拓。既然我们刚修正了该程序，对这个引理也必须有所交代：

记 $\{\widehat{W}_w\}$ 为 $M'$ 所对应的 $\widehat{V}$ 的滤过，$\{\mathcal{E}^{M'}_w\}$、$\{\mathcal{E}^M_w\}$ 为 $\mathcal{E}^{M'}$ 与 $\mathcal{E}^M$ 中相应的滤过。要证明的是：对一切 $w$ 有
\[
\deg \mathcal{E}^{M'}_w < \deg \mathcal{E}^M_w.
\]
这只须注意：由 \cite{lafforgue98} 第 2a 段命题 2，可以沿各子空间 $\widehat{W}_w$ 用一串（全部正向的）初等变换从 $M$ 得到 $M'$，因此对一切 $w$ 有
\[
\deg \mathcal{E}^{M'}_w < d - \deg \overline{\mathcal{E}}^{M'}_w \le \deg \mathcal{E}^M_w.
\]

\item 延拓唯一性的证明（要短得多），参看 \cite{lafforgue98} 一文的最后两页。
\end{itemize}
\end{remark*}

\origpage{56}
